% IV. Програмна реализация.
% Обща структура, йерархия на класовете, коментар на изходния текст.
% Разделът се попълва успоредно с писането на кода. Дотогава описва замисъла,
% а не постигнатото: твърдения за готово поведение тук не се пишат предварително.
\section{Програмна реализация}
\label{sec:implementation}

\subsection{Обща структура на изходния текст}

Проектът е организиран в три дървета. В \code{include/} стоят публичните заглавни
файлове, тоест единственото, което вижда приложението, използващо библиотеката. В
\code{src/} стои реализацията, разделена на подпапки по слоевете от
Раздел~\ref{sec:design}. В \code{tests/} стоят тестовете, писани с GoogleTest.
Разделянето на публично от вътрешно е съзнателно: то позволява представянето на
индексите да се промени, без да се променя нито един ред в кода на потребителя на
библиотеката.

Изграждането е с CMake. Библиотеката се изгражда като статична по подразбиране, тъй
като целевият случай на употреба е вграждане. Тестовете са отделна цел, която може
да бъде изключена, за да не изисква изтегляне на GoogleTest при вграждане в чужд
проект.

\subsection{Йерархия на типовете}

Ключовите типове следват слоевете. Тип за термин, който представя IRI, литерал или
празен възел. Тип за идентификатор на термин, който е обвивка над цяло число и не
допуска неявно превръщане, за да не се разменят по грешка идентификатори от различни
речници. Тип за триплет от три идентификатора. Речник, който предоставя двете посоки
на съпоставянето. Индекс, който предоставя обхождане по префикс. Възли на
алгебричното дърво, обединени в един вариант, вместо в йерархия с виртуални методи,
защото над дървото се извършват няколко различни обхождания, а не едно.

\TODO{диаграма на класовете и на зависимостите между модулите}

\subsection{Синтактични анализатори}

Двата анализатора споделят общ четец на знаци, който води ред и колона и предоставя
предварителен поглед с един знак. Върху него анализаторът на Turtle разпознава
префиксните обявявания, съкратените форми и вложените описания, а анализаторът на
SPARQL разпознава конструкциите на езика. Общият четец е причината съобщенията за
грешка да имат еднакъв вид и в двата случая.

\TODO{коментар на изходния текст на анализатора на SPARQL с листинг на
представителна функция, след като бъде написан}

\subsection{Механизъм за изпълнение}

Изпълнението е организирано като верига от итератори, всеки от които произвежда
следващото решение при поискване. Този модел е избран, за да не се материализират
междинните резултати изцяло и за да може ограничението на броя резултати да спре
работата рано. Съединяването на два подредени входа се извършва чрез сливане, а при
несъвместими подредби се вмъква стъпка за подреждане.

\TODO{описание на реализираните оператори и на тези, които остават извън обхвата}
